\unirsection{A fondo}

Los siguientes libros pueden servir de material de apoyo y para profundizar más en los contenidos de este tema.

\textbf{Yanofsky, N. S. y Mannucci, M. A. Quantum Computing for Computer Scientists. Cambridge University Press, 2019.}

El capítulo 1 introduce los números complejos directamente en el contexto de la mecánica cuántica, enfatizando las amplitudes y su papel probabilístico. Es la referencia más directamente conectada con la motivación cuántica de este tema.

\textbf{O'Neill, P. V. Matemáticas avanzadas para Ingeniería. Cengage Learning, 2012.}

El capítulo 8 cubre el álgebra y la geometría de los números complejos con abundantes ejemplos y ejercicios. Especialmente recomendable para afianzar el cálculo en forma polar y exponencial, y para el estudio de las funciones complejas elementales.

\textbf{Asmar, N. K. y Grafakos, L. Complex Analysis with Applications. Springer, 2018.}

El capítulo 1 proporciona una exposición rigurosa de los números complejos, incluyendo los aspectos topológicos del plano complejo que en este tema se tratan de forma más intuitiva.

\textbf{Lang, S. Complex Analysis. Springer, 2003.}

El capítulo 1 ofrece un tratamiento algebraico preciso, con especial atención a las propiedades del módulo y el conjugado y a la geometría del plano complejo. Recomendado para quienes deseen una visión más formal.